\subsection*{CU-29: Enviar un emote}
\addcontentsline{toc}{subsection}{CU-29: Enviar un emote}
\label{cu-29}

\begin{fichacu}{CU-29}{Enviar un emote}
  \crudFila{Responsable}{Ángel Gabriel Aguilar Hernández}
  \crudFila{Actor principal}{Jugador}
  \crudFila{Actores de sistema}{Servidor de partidas}
  \crudFila{Prototipos}{12e, 1a, 8a}
  \crudFila{Objetivo}{Permitir que el Jugador exprese una reacción rápida durante la partida con uno de los seis emotes de su conjunto equipado, sin escribir y sin interrumpir el juego.}
  \crudFila{Disparador}{El Jugador pulsa uno de los atajos del 1 al 6, o selecciona un emote en el selector de emotes de la \texttt{GUI\_\allowbreak{}Match}.}
  \textbf{Precondiciones} & \cuCelda{\begin{cureglas}
      \item \textbf{\mbox{PRE-01}} El Jugador tiene una \textit{Participacion} sin terminar en una \textit{Partida} en curso que no es contra la IA.
      \item \textbf{\mbox{PRE-02}} El Jugador tiene un objeto equipado en la \textit{Ranura} de emotes, que contiene seis emotes.
      \item \textbf{\mbox{PRE-03}} El Servidor de partidas está en ejecución.
    \end{cureglas}} \\ \hline
  \textbf{Flujo normal} & \cuCelda{\begin{cuenum}[start=1]
      \item El Jugador abre el selector de emotes o pulsa directamente un atajo del 1 al 6. \textit{(Ver \mbox{FA-01}, \mbox{FA-07})}
      \item El SJM verifica localmente que hayan pasado al menos diez segundos desde el último emote enviado (\mbox{RN-02}). \textit{(Ver \mbox{FA-02})}
      \item El SJM muestra el emote sobre el propio avatar de inmediato y envía el mensaje \texttt{EMOTE\_\allowbreak{}REQUEST} con el número de emote. \textit{(Ver \mbox{EX-01}, \mbox{EX-02})}
      \item El Servidor de partidas verifica que la conexión corresponda a una \textit{Participacion} sin terminar de esa partida. \textit{(Ver \mbox{FA-03})}
      \item El Servidor de partidas vuelve a verificar el intervalo de diez segundos con su propio registro en memoria (\mbox{RN-02}). \textit{(Ver \mbox{FA-04})}
      \item El Servidor de partidas toma el conjunto de emotes de la \textit{ParticipacionObjeto} de la ranura de emotes del Jugador, fijado al empezar la partida, y verifica que el número recibido exista en él. \textit{(Ver \mbox{FA-05})}
      \item El Servidor de partidas reenvía el emote a los demás participantes con \texttt{EMOTE\_\allowbreak{}RECEIVED}, sin escribir nada en la base de datos (\mbox{D-10}).
      \item El SJM de cada rival muestra el emote junto al avatar del Jugador durante unos segundos, salvo que ese rival haya silenciado los emotes del Jugador. \textit{(Ver \mbox{FA-06})}
    \end{cuenum}} \\
   & \cuCelda{\begin{cuenum}[start=9]
      \item Fin del caso de uso.
    \end{cuenum}} \\ \hline
  \textbf{Flujos alternos} & \cuCelda{\textbf{\mbox{FA-01}: El Jugador abre el selector completo}\par
    \begin{cuenum}[start=1]
      \item El Jugador selecciona el icono de emotes.
      \item El SJM despliega el selector con los seis emotes de su conjunto y el atajo de cada uno.
      \item El Jugador elige uno y el flujo continúa en el paso 2 del FN.
    \end{cuenum}} \\
   & \cuCelda{\textbf{\mbox{FA-02}: No han pasado diez segundos desde el último emote}\par
    \begin{cuenum}[start=1]
      \item El SJM no envía nada y muestra sobre el selector el tiempo que falta, sin ventanas.
      \item Fin del caso de uso.
    \end{cuenum}} \\
   & \cuCelda{\textbf{\mbox{FA-03}: El Jugador ya no participa en la partida}\par
    \begin{cuenum}[start=1]
      \item El Servidor de partidas descarta el emote sin reenviarlo: la partida terminó o el Jugador abandonó.
      \item Fin del caso de uso.
    \end{cuenum}} \\
   & \cuCelda{\textbf{\mbox{FA-04}: El servidor detecta emotes demasiado seguidos}\par
    \begin{cuenum}[start=1]
      \item El Servidor de partidas descarta el emote y responde con \texttt{EMOTE\_\allowbreak{}THROTTLED}.
      \item Si el Jugador supera el límite de forma repetida, el Servidor de partidas lo anota en la bitácora del servidor como indicio de cliente modificado.
      \item Fin del caso de uso.
    \end{cuenum}} \\
   & \cuCelda{\textbf{\mbox{FA-05}: El emote no pertenece al conjunto equipado}\par
    \begin{cuenum}[start=1]
      \item El Servidor de partidas descarta el emote, porque el aspecto del jugador se fijó al empezar y no puede cambiarse a mitad de partida (\mbox{RN-03} del \mbox{CU-14}).
      \item El Servidor de partidas lo anota en la bitácora del servidor como indicio de cliente modificado.
      \item Fin del caso de uso.
    \end{cuenum}} \\
   & \cuCelda{\textbf{\mbox{FA-06}: Un rival silencia los emotes del Jugador}\par
    \begin{cuenum}[start=1]
      \item El rival selecciona ``silenciar emotes'' sobre el avatar del Jugador.
      \item El SJM del rival deja de mostrar los emotes de ese jugador durante el resto de la partida, sin avisar al Jugador y sin enviar nada al servidor (\mbox{RN-04}).
      \item Fin del caso de uso.
    \end{cuenum}} \\
   & \cuCelda{\textbf{\mbox{FA-07}: El Jugador es espectador}\par
    \begin{cuenum}[start=1]
      \item El SJM no muestra el selector de emotes a los espectadores.
      \item Fin del caso de uso.
    \end{cuenum}} \\ \hline
  \textbf{Excepciones} & \cuCelda{\textbf{\mbox{EX-01}: Se agota el tiempo de espera o la conexión se interrumpe}\par
    \begin{cuenum}[start=1]
      \item El SJM no reenvía el emote: es efímero, y llegar tarde no tiene sentido.
      \item Si la conexión se perdió, el flujo continúa en el \mbox{CU-26}.
      \item Fin del caso de uso.
    \end{cuenum}} \\
   & \cuCelda{\textbf{\mbox{EX-02}: El mensaje está malformado}\par
    \begin{cuenum}[start=1]
      \item El Servidor de partidas descarta el mensaje sin responder.
      \item Fin del caso de uso.
    \end{cuenum}} \\ \hline
  \textbf{Postcondiciones} & \cuCelda{\begin{cureglas}
      \item \textbf{\mbox{POST-01}} Si el caso termina por el FN, los rivales que no silenciaron al Jugador vieron el emote.
      \item \textbf{\mbox{POST-02}} Ninguna estructura de la base de datos se modificó: el emote no deja rastro persistente.
      \item \textbf{\mbox{POST-03}} El Jugador no puede enviar otro emote durante los diez segundos siguientes.
    \end{cureglas}} \\ \hline
  \textbf{Reglas de negocio} & \cuCelda{\begin{cureglas}
      \item \textbf{\mbox{RN-01}} En una partida hay seis emotes disponibles, con atajos del 1 al 6, que son los del conjunto equipado en la \textit{Ranura} de emotes.
      \item \textbf{\mbox{RN-02}} Un jugador puede enviar a lo sumo un emote cada diez segundos. El cliente lo aplica para responder al instante y el servidor lo repite para decidir.
      \item \textbf{\mbox{RN-03}} Los emotes no se persisten (\mbox{D-10}). Se reenvían y se pintan.
      \item \textbf{\mbox{RN-04}} Silenciar los emotes de un rival es local a la partida y al cliente: no se guarda y no se comunica al rival. Es distinto de silenciar a un jugador en el \mbox{CU-33}, que afecta también al chat y se conserva.
      \item \textbf{\mbox{RN-05}} No hay emotes en partidas contra la IA ni para los espectadores.
    \end{cureglas}} \\ \hline
  \textbf{Observación} & \cuCelda{\textit{Sobre por qué los emotes no se guardan y los mensajes sí.}\par
    El chat se guarda porque el \mbox{CU-42} lo adjunta como prueba cuando alguien reporta. Los emotes no pueden ofender del mismo modo: son seis imágenes elegidas por el propio sistema, sin texto libre. El único abuso posible es la insistencia, y ese ya lo impiden el intervalo de diez segundos y la opción de silenciarlos. Guardarlos añadiría una escritura cada pocos segundos por jugador sin que ninguna funcionalidad la lea.} \\ \hline
  \textbf{Datos que toca} & \cuCelda{\begin{cudatos}
      \item \textbf{\textit{Participacion}:} consulta en memoria del servidor, sin acceso a la base de datos \textit{(paso 4)}
      \item \textbf{\textit{ParticipacionObjeto}:} consulta en memoria el conjunto de emotes fijado al empezar \textit{(paso 6)}
      \item \textbf{—:} ninguna escritura en la base de datos
    \end{cudatos}} \\ \hline
\end{fichacu}
\clearpage
